\documentclass[12pt, a4paper]{ctexart}
\usepackage{fontspec}
\usepackage{xcolor}
\usepackage{hyperref}
\usepackage{geometry}
\usepackage{enumitem}
\usepackage{parskip}
\usepackage{titlesec}
\usepackage{tcolorbox}
\tcbuselibrary{breakable}
\tcbset{autoparskip}

\usepackage{setspace}
\usepackage{listings}
\usepackage{amssymb}

\geometry{left=2.5cm, right=2.5cm, top=2.5cm, bottom=2.5cm}

\setmainfont{Times New Roman}
\setCJKmainfont{SimSun}

\hypersetup{
    colorlinks=true,
    linkcolor=blue!70!black,
    urlcolor=cyan,
    pdftitle={多线程-逐题精讲解义},
    pdfauthor={专升本-fifine老师}
}

\definecolor{myblue}{RGB}{30,100,200}
\definecolor{lightblue}{RGB}{230,240,255}
\definecolor{darkblue}{RGB}{0,50,120}

\newtcolorbox{bluebox}[1][]{
    breakable,
    colback=lightblue,
    colframe=myblue,
    coltitle=darkblue,
    fonttitle=\bfseries,
    title={#1},
    boxrule=0.8pt,
    arc=4pt,
    left=6pt,
    right=6pt,
    top=6pt,
    bottom=6pt,
    before skip=8pt,
    after skip=8pt
}

\usepackage{etoolbox}
\BeforeBeginEnvironment{bluebox}{\par\vfill}%
\AfterEndEnvironment{bluebox}{\par\vfill}%

\titleformat{\section}
  {\normalfont\Large\bfseries\color{myblue}}{\thesection}{1em}{}
\titlespacing*{\section}{0pt}{3.5ex plus 1ex minus .2ex}{1.5ex plus .2ex}

\title{\textcolor{myblue}{\textbf{多线程 - 逐题精讲解义}}}
\author{专升本-fifine老师 教学组}
\date{2026年03月05日}

\lstset{
    language=Java,
    basicstyle=\ttfamily\small,
    keywordstyle=\color{blue},
    commentstyle=\color{green!60!black},
    stringstyle=\color{orange},
    numbers=left,
    numberstyle=\tiny\color{gray},
    frame=single,
    breaklines=true,
    columns=fullflexible,
    showstringspaces=false
}

\newcommand{\code}[1]{\texttt{#1}}

\begin{document}

\maketitle
\thispagestyle{empty}
\newpage

\section{多线程}

\begin{enumerate}[label=\textbf{\arabic*.}]

	\item \textbf{以下关于线程的说法错误的是}
	      \begin{enumerate}[label=\Alph*.]
		      \item 多线程可以提高程序的执行效率
		      \item 线程之间共享内存
		      \item 实现线程的方式有继承Thread类和实现Runnable接口
		      \item 线程是进程的一部分
	      \end{enumerate}

	      \begin{bluebox}{答案与深度解析}
		      \textbf{答案：B. 线程之间共享内存}

		      \textbf{【核心认知框架】}
		      \begin{itemize}[label={}, leftmargin=*, nosep]
			      \item \textbf{题目本质}：考查Java线程的内存模型
			      \item \textbf{关键区分点}：共享堆内存 vs 独立栈内存
			      \item \textbf{命题人意图}：测试对线程内存分配机制的准确理解
		      \end{itemize}

		      \textbf{【选项原理深度拆解】}

		      \begin{itemize}[leftmargin=*, nosep, label={}]
			      \item[A] \textbf{多线程可以提高程序的执行效率 — 正确}
			            \begin{itemize}[label={}, leftmargin=*, nosep]
				            \item \textbf{适用场景}：
				                  \begin{itemize}[label={}, nosep]
					                  \item CPU密集型任务：多核并行计算
					                  \item IO密集型任务：线程阻塞时其他线程继续执行
					                  \item 异步任务：主线程不阻塞
				                  \end{itemize}
				            \item \textbf{效率提升原因}：
				                  \begin{itemize}[label={}, nosep]
					                  \item CPU利用率提升（多核）
					                  \item IO等待时CPU可处理其他任务
					                  \item 响应速度快（UI线程不阻塞）
				                  \end{itemize}
				            \item \textbf{代码示例}：
				                  \begin{lstlisting}[language=Java]
// 单线程执行
long start = System.currentTimeMillis();
for (int i = 0; i < 10; i++) {
    processTask(i);
}
long singleTime = System.currentTimeMillis() - start;

// 多线程并行执行
start = System.currentTimeMillis();
ExecutorService executor = Executors.newFixedThreadPool(10);
List<Future<?>> futures = new ArrayList<>();
for (int i = 0; i < 10; i++) {
    final int index = i;
    futures.add(executor.submit(() -> processTask(index)));
}
for (Future<?> future : futures) {
    future.get();
}
executor.shutdown();
long multiTime = System.currentTimeMillis() - start;

// 多线程执行时间通常比单线程短
// 在多核CPU上，理论上可达N倍加速（N=核心数）
				      \end{lstlisting}
			            \end{itemize}

			      \item[B] \textbf{线程之间共享内存 — 错误（本题答案）}
			            \begin{itemize}[label={}, leftmargin=*, nosep]
				            \item \textbf{错误原因}：表述不准确，应该是"部分共享"
				            \item \textbf{内存实际分配}：
				                  \begin{itemize}[label={}, nosep]
					                  \item 堆内存（Heap）：所有线程共享
					                  \item 方法区（Method Area）：所有线程共享
					                  \item 栈内存（Stack）：每个线程独立
					                  \item 程序计数器（PC）：每个线程独立
					                  \item 本地方法栈：每个线程独立
				                  \end{itemize}
				            \item \textbf{线程内存模型图}：
				                  \begin{lstlisting}[language=Java]
Java内存模型（JMM）：
┌─────────────────────────────────────┐
│     方法区（堆内存）- 所有线程共享    │
│  ┌───────┐  ┌───────┐  ┌───────┐   │
│  │ 对象1 │  │ 对象2 │  │ 对象3 │   │
│  └───────┘  └───────┘  └───────┘   │
└─────────────────────────────────────┘
        ↑            ↑            ↑
   线程1共享   线程2共享   线程3共享
        │            │            │
┌───────┴────┐ ┌─────┴─────┐ ┌────┴────┐
│ 线程1栈    │ │ 线程2栈   │ │ 线程3栈  │
│ (独立私有) │ │ (独立私有) │ │ (独立私有)│
└────────────┘ └───────────┘ └─────────┘
// 堆内存共享，栈内存独立
					      \end{lstlisting}
				            \item \textbf{代码示例}：
				                  \begin{lstlisting}[language=Java]
public class MemorySharing {
    // 静态变量（堆内存）：所有线程共享
    static int sharedCounter = 0;

    public static void main(String[] args) {
        SharedObject obj = new SharedObject();  // 对象在堆中

        // 线程1访问共享对象
        new Thread(() -> {
            obj.increment();  // 可见线程2的修改
            System.out.println("Thread1: " + sharedCounter);
        }).start();

        // 线程2访问共享对象
        new Thread(() -> {
            obj.increment();
            System.out.println("Thread2: " + sharedCounter);
        }).start();

        // 本地变量（栈内存）：每个线程独立
        Thread t1 = new Thread(() -> {
            int localVar = 100;  // 线程1栈中的变量
        });
        Thread t2 = new Thread(() -> {
            int localVar = 200;  // 线程2栈中的变量
            // localVar1和localVar2互不影响
        });
    }
}

class SharedObject {
    // 对象字段在堆中，多线程共享
    private int counter = 0;

    public void increment() {
        this.counter++;
    }
}
// 线程共享堆内存中的对象和静态变量
// 但每个线程有独立的栈空间存放局部变量
				      \end{lstlisting}
				            \item \textbf{重点提示}：
				                  \begin{itemize}[label={}, nosep]
					                  \item 可以共享：堆中的实例对象、静态变量、数组
					                  \item 不共享：局部变量、方法参数、异常表、栈帧
					                  \item 同一块栈内存不能被多个线程同时访问（线程安全）
				                  \end{itemize}
			            \end{itemize}

			      \item[C] \textbf{实现线程的方式有继承Thread类和实现Runnable接口 — 正确}
			            \begin{itemize}[label={}, leftmargin=*, nosep]
				            \item \textbf{方式1：继承Thread类}
				                  \begin{lstlisting}[language=Java]
class MyThread extends Thread {
    @Override
    public void run() {
        System.out.println("线程运行中");
    }
}

// 使用
MyThread thread = new MyThread();
thread.start();  // 启动新线程
// 缺点：Java不支持多继承
				            \end{lstlisting}
				            \item \textbf{方式2：实现Runnable接口}
				                  \begin{lstlisting}[language=Java]
class MyRunnable implements Runnable {
    @Override
    public void run() {
        System.out.println("线程运行中");
    }
}

// 使用
Thread thread = new Thread(new MyRunnable());
thread.start();
// 优点：可继承其他类，推荐使用
				            \end{lstlisting}
				            \item \textbf{方式3：实现Callable接口（可返回结果）}
				                  \begin{lstlisting}[language=Java]
class MyCallable implements Callable<Integer> {
    @Override
    public Integer call() throws Exception {
        return 42;
    }
}

// 使用
FutureTask<Integer> future = new FutureTask<>(new MyCallable());
new Thread(future).start();
Integer result = future.get();  // 获取返回值
// Callable支持返回值和抛出异常
				            \end{lstlisting}
			            \end{itemize}

			      \item[D] \textbf{线程是进程的一部分 — 正确}
			            \begin{itemize}[label={}, leftmargin=*, nosep]
				            \item \textbf{进程vs线程}：
				                  \begin{itemize}[label={}, nosep]
					                  \item 进程：资源分配的基本单位（独立的内存空间）
					                  \item 线程：CPU调度的基本单位（共享进程资源）
					                  \item 一个进程包含多个线程
					                  \item 一个进程至少有一个main线程
				                  \end{itemize}
				            \item \textbf{内存共享关系}：
				                  \begin{lstlisting}[language=Java]
进程资源分配：
┌─────────────────────────┐
│     进程                │
│  ┌───────────────────┐  │
│  │  代码段（只读）    │  │ ← 所有线程共享
│  │  数据段            │  │ ← 所有线程共享
│  │  堆内存            │  │ ← 所有线程共享
│  │  ┌─ 线程1 栈 ┐    │  │ ← 线程1私有
│  │  │  局部变量 │    │  │
│  │  └───────────┘    │  │
│  │  ┌─ 线程2 栈 ┐    │  │ ← 线程2私有
│  │  │  局部变量 │    │  │
│  │  └───────────┘    │  │
│  └───────────────────┘  │
└─────────────────────────┘
// 线程是进程内的执行单元
					      \end{lstlisting}
				            \item \textbf{主线程}：JVM启动时创建的main()执行线程
			            \end{itemize}
		      \end{itemize}

		      \textbf{【应背诵知识点】}
		      \begin{itemize}[label={}, nosep]
			      \item 多线程提高执行效率（有条件）
			      \item 线程共享堆和方法区，每个线程有独立栈
			      \item 线程创建方式：继承Thread、实现Runnable、实现Callable
			      \item 线程是进程内的执行单元，共享进程资源
			      \item 口诀：线程共享堆内存，独立栈不共享，提高效率条件多，进程包含若干线
		      \end{itemize}

		      \textbf{【命题趋势与实战策略】}
		      \textbf{考场秒杀}：看到"线程之间共享内存"→注意表述，堆共享栈独立
		      \textbf{记忆策略}：堆是共享仓库，栈是私有工作台
	      \end{bluebox}

\end{enumerate}

\end{document}
